\documentclass[10pt,a4paper]{amsart}
\usepackage{amsmath,amssymb,amsthm,graphicx}
% Retained minimal preamble for the Stacks Project chapter varieties.tex.
% Verbatim excerpt of the chapter preamble: Xy-pic setup (native preamble lines 14-18), the native theorem
% declarations with the mathematical macros (native preamble lines 154-217), and the project own no-op
% definitions of the reference/slogan/history display environments (native wrapper.tex lines 1-3) that the
% counted statement uses for its slogan. Copied mechanically; nothing added, removed or redefined.
% For commutative diagrams we use Xy-pic
\usepackage[all]{xy}

% We use 2cell for 2-commutative diagrams.
\xyoption{2cell}
% Theorem environments.
%
\theoremstyle{plain}
\newtheorem{theorem}[subsection]{Theorem}
\newtheorem{proposition}[subsection]{Proposition}
\newtheorem{lemma}[subsection]{Lemma}

\theoremstyle{definition}
\newtheorem{definition}[subsection]{Definition}
\newtheorem{example}[subsection]{Example}
\newtheorem{exercise}[subsection]{Exercise}
\newtheorem{situation}[subsection]{Situation}

\theoremstyle{remark}
\newtheorem{remark}[subsection]{Remark}
\newtheorem{remarks}[subsection]{Remarks}

\numberwithin{equation}{subsection}

% Macros
%
\def\lim{\mathop{\mathrm{lim}}\nolimits}
\def\colim{\mathop{\mathrm{colim}}\nolimits}
\def\Spec{\mathop{\mathrm{Spec}}}
\def\Hom{\mathop{\mathrm{Hom}}\nolimits}
\def\Ext{\mathop{\mathrm{Ext}}\nolimits}
\def\SheafHom{\mathop{\mathcal{H}\!\mathit{om}}\nolimits}
\def\SheafExt{\mathop{\mathcal{E}\!\mathit{xt}}\nolimits}
\def\Sch{\mathit{Sch}}
\def\Mor{\mathop{\mathrm{Mor}}\nolimits}
\def\Ob{\mathop{\mathrm{Ob}}\nolimits}
\def\Sh{\mathop{\mathit{Sh}}\nolimits}
\def\NL{\mathop{N\!L}\nolimits}
\def\CH{\mathop{\mathrm{CH}}\nolimits}
\def\proetale{{pro\text{-}\acute{e}tale}}
\def\etale{{\acute{e}tale}}
\def\QCoh{\mathit{QCoh}}
\def\Ker{\mathop{\mathrm{Ker}}}
\def\Im{\mathop{\mathrm{Im}}}
\def\Coker{\mathop{\mathrm{Coker}}}
\def\Coim{\mathop{\mathrm{Coim}}}

% Boxtimes
%
\DeclareMathSymbol{\boxtimes}{\mathbin}{AMSa}{"02}

%
% Macros for moduli stacks/spaces
%
\def\QCohstack{\mathcal{QC}\!\mathit{oh}}
\def\Cohstack{\mathcal{C}\!\mathit{oh}}
\def\Spacesstack{\mathcal{S}\!\mathit{paces}}
\def\Quotfunctor{\mathrm{Quot}}
\def\Hilbfunctor{\mathrm{Hilb}}
\def\Curvesstack{\mathcal{C}\!\mathit{urves}}
\def\Polarizedstack{\mathcal{P}\!\mathit{olarized}}
\def\Complexesstack{\mathcal{C}\!\mathit{omplexes}}
% \Pic is the operator that assigns to X its picard group, usage \Pic(X)
% \Picardstack_{X/B} denotes the Picard stack of X over B
% \Picardfunctor_{X/B} denotes the Picard functor of X over B
\def\Pic{\mathop{\mathrm{Pic}}\nolimits}
\def\Picardstack{\mathcal{P}\!\mathit{ic}}
\def\Picardfunctor{\mathrm{Pic}}
\def\Deformationcategory{\mathcal{D}\!\mathit{ef}}
\newenvironment{reference}{}{}
\newenvironment{slogan}{}{}
\newenvironment{history}{}{}
\usepackage{cleveref}
\crefname{originaldefinition}{original definition}{original definitions}
\Crefname{originaldefinition}{Original definition}{Original definitions}
\makeatletter
\newlabel{algebra-definition-geometrically-regular}{{0382}{source}{}{sourcepointer-1}{}}
\newlabel{algebra-definition-geometrically-regular@cref}{{[originaldefinition][0][]0382}{[1][0][]source}{}{}{}}
\makeatother
\crefname{originallemma}{original lemma}{original lemmas}
\Crefname{originallemma}{Original lemma}{Original lemmas}
\makeatletter
\newlabel{algebra-lemma-characterize-local-ring}{{00E9}{source}{}{sourcepointer-2}{}}
\newlabel{algebra-lemma-characterize-local-ring@cref}{{[originallemma][0][]00E9}{[1][0][]source}{}{}{}}
\makeatother
\crefname{originallemma}{original lemma}{original lemmas}
\Crefname{originallemma}{Original lemma}{Original lemmas}
\makeatletter
\newlabel{algebra-lemma-p-ring-map}{{0BRA}{source}{}{sourcepointer-3}{}}
\newlabel{algebra-lemma-p-ring-map@cref}{{[originallemma][0][]0BRA}{[1][0][]source}{}{}{}}
\makeatother
\crefname{originallemma}{original lemma}{original lemmas}
\Crefname{originallemma}{Original lemma}{Original lemmas}
\makeatletter
\newlabel{morphisms-lemma-base-change-surjective}{{01S1}{source}{}{sourcepointer-4}{}}
\newlabel{morphisms-lemma-base-change-surjective@cref}{{[originallemma][0][]01S1}{[1][0][]source}{}{}{}}
\makeatother
\title{Pointwise characterizations of geometric regularity by purely inseparable base change and by the local ring}
\author{The Stacks Project Authors}
\date{stacks-project:varieties.tex:local-snapshot-20261011; GNU FDL 1.2 or later}
\begin{document}
\maketitle
\noindent\textit{Editorial note.} All delivered mathematical text is the author's own, copied from the pinned Stacks Project chapter file varieties.tex; no statement or proof step was rewritten, reordered or removed.\par\medskip
\noindent\textit{Editorial note.} The proof is complete as written and closes the cycle of equivalences; no step is referred to the reader.\par\medskip
\noindent\textit{Editorial note.} The counted spans carry 4 navigation stubs and no bibliography citation.\par\medskip
\noindent\textit{Editorial note.} The preamble is the chapter's retained minimal preamble (sources/stacks-varieties/retained-preamble.tex lines 1-77), a verbatim excerpt of the native preamble.\par\medskip
\noindent\textit{Editorial note.} The label lemma-geometrically-regular-at-point is defined only in varieties.tex, so no cross-chapter restatement had to be weighed for this card.\par\medskip
\noindent\textit{Allowed standard prerequisites (editorial).}\begin{itemize}
\item Geometric regularity of a k-algebra. Algebra, Definition algebra-definition-geometrically-regular, invoked at proof lines 2409-2410 and named in the statement. Supplies the algebraic property identified with geometric regularity of the local ring.
\item Spectrum of a purely inseparable base change and localness of the tensor product. Algebra, Lemmas algebra-lemma-p-ring-map and algebra-lemma-characterize-local-ring, invoked at proof lines 2403-2406. Supplies that the tensor product of the local ring with a finite purely inseparable extension is again local, with the same spectrum.
\item Surjectivity is stable under base change. Morphisms, Lemma morphisms-lemma-base-change-surjective, invoked at proof lines 2413-2416. Supplies the existence of a point above x after an arbitrary field extension.
\end{itemize}\medskip
\ifdefined\hypertarget\hypertarget{sourcepointer-1}{}\fi
\noindent\textit{Source reference.} Original definition 0382: sources/stacks-chow/tags.txt, source lines 4221--4221. Pointer outside the counted span, invoked once for context; the load-bearing argument of the counted statement is the author's proof reproduced in full in the proof span.\par\medskip
\ifdefined\hypertarget\hypertarget{sourcepointer-2}{}\fi
\noindent\textit{Source reference.} Original lemma 00E9: sources/stacks-chow/tags.txt, source lines 526--526. Pointer outside the counted span, invoked once for context; the load-bearing argument of the counted statement is the author's proof reproduced in full in the proof span.\par\medskip
\ifdefined\hypertarget\hypertarget{sourcepointer-3}{}\fi
\noindent\textit{Source reference.} Original lemma 0BRA: sources/stacks-chow/tags.txt, source lines 14763--14763. Pointer outside the counted span, invoked once for context; the load-bearing argument of the counted statement is the author's proof reproduced in full in the proof span.\par\medskip
\ifdefined\hypertarget\hypertarget{sourcepointer-4}{}\fi
\noindent\textit{Source reference.} Original lemma 01S1: sources/stacks-chow/tags.txt, source lines 2322--2322. Pointer outside the counted span, invoked once for context; the load-bearing argument of the counted statement is the author's proof reproduced in full in the proof span.\par\medskip
\begin{lemma}
\label{lemma-geometrically-regular-at-point}
Let $k$ be a field.
Let $X$ be a locally Noetherian scheme over $k$.
Let $x \in X$.
The following are equivalent
\begin{enumerate}
\item $X$ is geometrically regular at $x$,
\item for every finite purely inseparable field extension $k'$ of $k$
and $x' \in X_{k'}$ lying over $x$ the local ring
$\mathcal{O}_{X_{k'}, x'}$ is regular, and
\item the ring $\mathcal{O}_{X, x}$ is geometrically
regular over $k$ (see
Algebra, Definition \ref{algebra-definition-geometrically-regular}).
\end{enumerate}
\end{lemma}

\begin{proof}
It is clear that (1) implies (2).
Assume (2). This in particular implies that $\mathcal{O}_{X, x}$
is a regular local ring. Let $k'/k$ be a finite purely inseparable
field extension. Consider the ring $\mathcal{O}_{X, x} \otimes_k k'$.
By Algebra, Lemma \ref{algebra-lemma-p-ring-map}
its spectrum is the same as the spectrum of $\mathcal{O}_{X, x}$.
Hence it is a local ring also
(Algebra, Lemma \ref{algebra-lemma-characterize-local-ring}).
Therefore there is a unique point $x' \in X_{k'}$ lying over $x$
and $\mathcal{O}_{X_{k'}, x'} \cong \mathcal{O}_{X, x} \otimes_k k'$.
By assumption this is a regular ring. Hence we deduce (3)
from the definition of a geometrically regular ring.

\medskip\noindent
Assume (3). Let $k'/k$ be a field extension. Since
$\Spec(k') \to \Spec(k)$ is surjective, also
$X_{k'} \to X$ is surjective
(Morphisms, Lemma \ref{morphisms-lemma-base-change-surjective}).
Let $x' \in X_{k'}$ be any point lying over $x$.
The local ring $\mathcal{O}_{X_{k'}, x'}$
is a localization of the ring $\mathcal{O}_{X, x} \otimes_k k'$.
Hence it is regular by assumption and (1) is proved.
\end{proof}

\end{document}
